\scanpage{84}\mspage{28}%
% L1-2
valence, par définition des Hot-fibrés, donc c'est une $W$-équivalence.

% L3-4
On a donc les foncteurs
\diagram{\begin{tikzcd}[ampersand replacement=\&,column sep=large]
\deleted{\unsure{\mathrm{Cat}_{/S}\longrightarrow}}\ \underline{\mathrm{Fibsc}}_0\,S
  \arrow[r,shift left,"p"] \& \mathrm{Cat}_{/S}\arrow[l,shift left,"\Phi_0"],
\end{tikzcd}}
% L5-8
il faut voir qu'ils satisfont aux conditions équivalentes du lemme p.~11. (et ce sera pareil pour \deleted{\unsure{p}}, $q,\Psi_0$). On sait déjà que $p$ est compatible avec les localiseurs, \deleted{et $\Phi_0$}
% L9-12
\deleted{aussi car il suffit de le voir pour le foncteur \unsure{(1)} dans $\underline{\mathrm{Fibsc}}\,S$ (les localiseurs s'induisent), et alors $\Phi_0=\Phi$} Prouvons le pour $\deleted{\unsure{\Psi|\Phi}}\unsure{\Phi|\Psi}_0$. Soit
% L13
\[
 f:X\unsure{\mathord{-}|\longrightarrow}Y
\]
% L14-15
dans $\mathrm{Cat}_{/S}$ une $W$-équivalence, alors $rC(f)$ l'est aussi. Il suffit de prouver ceci :

% L16
Soit \quad $g:X'\longrightarrow Y'$
% L17-20
un morphisme dans \deleted{\unsure{Hot}}$\underline{\mathrm{Fibhot}}\,S$ (p. exemple un morphisme de $\underline{\mathrm{Parf}}\,S$), qui soit une $W$-équivalence donc des $W$-équivalences sur les fibres. Alors $\deleted{\unsure{\Psi}}\Phi(g):\deleted{\unsure{\Psi}}\Phi(X')\unsure{\mathord{-}|\longrightarrow}\deleted{\unsure{\Psi}}\Phi(Y')$ est une $W$-équivalence fibre par fibre,
% L21-22
i.e. \quad $X'_{/\deleted{\unsure{S}}s}\longrightarrow Y'_{/s}$ une $W$-équiv. pour \deleted{\unsure{tout}} $s\in S$. Ça a été vu \unsure{dans|du} Cor.~2 \unsure{p.|page}~20.
% archive sheet number: 83
